\section{Reproducibility contract}\label{app:reproduction}

All commands below are run from the repository root with Python in UTF-8,
no-bytecode mode and the exact environment in \path{requirements.txt}.

\subsection{Manifest and core replay}

\begin{verbatim}
python -X utf8 -B scripts/verify.py --profile manifest
python -X utf8 -B scripts/verify.py --profile core
\end{verbatim}

The core profile checks the complete frozen manifest, manuscript/PDF receipt,
C32 and C46 realization certificates, the P21/P26/P49 specimen locks, both
residual-area producers, an integrated P21 smoke replay, archive-level
checksums, import closure, and repository policy.

\subsection{Full mutation and P21 replay}

\begin{verbatim}
python -X utf8 -B scripts/verify.py --profile full
\end{verbatim}

The full profile adds all three mutation suites.  Its integrated P21 verifier
reconstructs the 72 participating columns, 7,444
support-four relations, 1,119 target orbits, rank-nine quotient, minimum
triples, and Markov obstructions without importing the principal
support-four analysis file.  It then repeats the full replay in a temporary
tree containing only the manifest allowlist; this detects hidden dependencies
on virtual environments, caches, build debris, or unshipped files.

\subsection{Manuscript verification}

The manuscript static verifier checks the frozen section order, labels,
citations, load-bearing theorem fragments, no-firstness boundary, and
presence of the built PDF.  The PDF is then reopened, text-extracted for
sanity, rendered to PNG, and visually inspected page by page.  Build and
verification commands are in \texttt{paper/BUILD.md}.

The code and result filenames beginning with \texttt{coarea\_} are stable
legacy artifact identifiers.  In the mathematical prose the parameter is
called the largest-piece residual area.
